\documentclass[10pt,a4paper]{amsart}
\usepackage[margin=19mm]{geometry}
\usepackage{amsmath,amssymb,mathtools}
\usepackage[scr=rsfs,cal=euler]{mathalfa}
\usepackage{xfrac}
\usepackage[unicode,hidelinks,bookmarks=false]{hyperref}
\newcommand{\ie}{i.e.\ }
\newcommand{\cf}{cf.\ }
\newcommand{\resp}{respectively\ }
\newcommand{\Subsection}{\subsection}
\providecommand{\nobreakdash}{\nobreak\hbox{-}}
\setlength{\emergencystretch}{2em}
\multlinegap0pt
\theoremstyle{plain}
\newtheorem{thm}{Theorem}
\newtheorem{prop}[thm]{Proposition}
\newtheorem{lem}[thm]{Lemma}
\newtheorem{cor}[thm]{Corollary}
\theoremstyle{remark}
\newtheorem{remark}[thm]{Remark}
\DeclareMathOperator{\supp}{supp}
\newcommand{\C}{\mathbb{C}}
\newcommand{\D}{\mathbb{D}}
\newcommand{\isdef}{\overset{\mathrm{def}}{=\joinrel=}}
\title[Flow of the zeros of derivatives]{Flow of the zeros of polynomials under iterated differentiation: the difference of consecutive inverse branches up to order one over n squared}
\author{Andrei Mart\'{\i}nez-Finkelshtein, Evguenii A. Rakhmanov}
\date{Source: arXiv:2408.13851v3 (CC BY 4.0)}
\begin{document}
\maketitle
\noindent\emph{Source and licence.} Flow of the zeros of polynomials under iterated differentiation. Version arXiv:2408.13851v3 (CC BY 4.0), distributed under the Creative Commons Attribution 4.0 International licence, \url{http://creativecommons.org/licenses/by/4.0/}, as recorded in the arXiv licence metadata for this exact version and shown on the abstract page \url{https://arxiv.org/abs/2408.13851v3}. Full licence notice: licenses/arxiv-2408.13851-CC-BY-4.0.txt.\par\smallskip

\noindent\textit{Author byline.} The two authors are listed in the article's own byline spelling, which writes the accented form of the first family name; the arXiv metadata that accompanies this version uses unaccented transliterations of both names. No author is added or dropped.\par\smallskip

\noindent\textit{Editorial scope.} Counted unit: the article's Proposition with source label \texttt{lem5} (source0.tex lines 851--863), the expansion of the difference of two consecutive inverse branches of the differentiated-polynomial flow into the explicit singular term minus one over n zeta plus a remainder, together with the analyticity and the quadratic-in-one-over-n bound for that remainder on a smaller disc; delivered with the article's complete original proof of it (source0.tex lines 864--976, one \texttt{proof} environment of 113 lines, adjacent to the statement). No mathematics is written, repaired or re-derived: statement and proof are the article's own text line for line, in the article's own notation (the inverse branches, the discs, the Cauchy transform bounds, the auxiliary functions and their derivatives, the remainder terms), and no retained line is substituted, re-flowed or omitted. Required context retained uncounted, in source order: the labelled display bounding the inverse branches at infinity (lines 393--396); the labelled display of the change of variables for the Cauchy transform (lines 398--400); the labelled display of the normalisation condition on the measure (lines 473--476); and the statement of the article's Cauchy-transform estimate (lines 480--496), including its labelled two-sided bounds, which the counted proof invokes. The proofs of those context results are not delivered. Nothing else of the article is retained: its main theorems, its other propositions and its appendices are omitted. No citation command occurs in any retained line, so the article's bibliography is not delivered and the wrapper carries no bibliography entry. Every cross-reference inside the retained spans resolves inside the delivered document. Editorial formatting only: an AMS \texttt{amsart} preamble that restates the environments and macros the retained lines use, namely the article's declarations of Theorem, Proposition, Lemma, Corollary and Remark sharing one counter, the article's support operator and its three macros for the complex plane, the unit disc and the defining equality. The article's own source declares the class \texttt{amsart} as well; no class or style file is shipped with this package. Numbering differs from the article, because the article resets its shared theorem counter in each section and because only a subset of environments is retained: the retained Cauchy-transform estimate prints as Proposition 1 and the counted proposition as Proposition 2, whereas the article prints them in its own sections. No picture environment and no graphical inclusion occurs in any retained line, and the package delivers no image asset. One bundled source file, \texttt{sources/164249/source0.tex}, accompanies the delivered item as the exact source of every retained span. Every retained span is recorded in \texttt{delivery-evidence.json} with method \texttt{exact\_tex}.
\begin{equation}
    \label{estimateatinftyvnk}
v_{n,k} (z) = \frac{n-k}{n  }\, \frac{1}{  z } + \mathcal O(z^{-2}), \quad z\to \infty.
\end{equation}

\begin{equation} \label{eqCauchyTransfNew}
	v_{n,k+1}(z)- v_{n,k}(z) =\frac{1}{n}\, \frac{\partial_z\, v_{n,k}(z)}{  v_{n,k}(z)}. 
\end{equation}

\begin{equation}
    \label{normCondition}
m_0=\mu(\C)=1,
\end{equation}

\begin{prop}
    \label{prob:boundsCauchy}
    Let $\mu$ be a finite Borel and compactly-supported signed measure on $\C$, such that \eqref{normCondition} holds, and let  $0<r<+\infty$ be such that $\supp (\mu) \subset\D_r$. Then for any $R>  (1+|\mu|(\C))\, r$,  $\mathcal C^\mu$ is analytic in $\overline{\Omega_R}$ and does not vanish for $z\in \Omega_R$.

If in addition $\mu$ is a positive (probability) measure, then for $|z|\ge R>2r$,
    \begin{align}
        \label{boundsCauchy1}
    & \frac{2}{9}< \frac{1-r/R}{(1+r/R)^2} \le \left|  z\mathcal C^\mu(z)\right| \le \frac{R}{R-r} ,
    \\  
    & \frac{1-2r/R}{(1+r/R)^4} \le \left| z^2\, \frac{d}{dz}\mathcal C^\mu(z)\right| \le \frac{1}{(1-r/R)^2}<4.
\label{boundsCauchy2}
    \end{align}
Moreover, if $R>(\sqrt{2}+1)r$, then $C^\mu$ is univalent (injective) in  $\Omega_R$, and 
\begin{equation} \label{image}
\overline{\D_{1/(4R)}}\subset C^\mu\left( \overline{\Omega_R} \right)\subset \overline{\D_{1/(R-r)}}.
\end{equation}
\end{prop}

\begin{prop}
    \label{lem5}
For the inverse functions $v_{n, k}^{-1}(\zeta)$ and $v_{n, k+1}^{-1}(\zeta)$, with $\zeta \in \D_\rho$, we have
\begin{equation}
    \label{eq:12}
v_{n,k+1}^{-1}\left(\zeta\right)-v_{n, k}^{-1}(\zeta)=-\frac{1}{n \zeta}  +\Delta_{n, k}(\zeta).
\end{equation}
There is $\rho'<\rho$ and $M>0$, depending only on $r$, such that $\Delta_{n, k}$ are analytic in $\D_{\rho'}$ and 
\begin{equation}
    \label{boundDelta}
    \left|\Delta_{n, k}(\zeta)\right| \leq \frac{M}{n^2 }, \quad \zeta \in \D_{\rho'}.
\end{equation}
\end{prop}

\begin{proof}
Notice that by \eqref{estimateatinftyvnk}, 
$$
v_{n,k+1}^{-1}\left(\zeta\right)-v_{n, k}^{-1}(\zeta) +\frac{1}{n \zeta} 
$$
is analytic at $\zeta=0$, so it is enough to establish the bound \eqref{boundDelta} on a circle $|\zeta|=\rho'$.

To simplify the notation, denote $\varphi_{n, k}(\zeta)\isdef v_{n, k}^{-1}(\zeta)$ and let $^\prime$ stand for the derivative with respect to the corresponding variable that should be clear in each context: $v_{n, k}'(z)=\partial_z \,   v_{n, k}(z)$, $\varphi_{n, k}'(\zeta)=\partial_\zeta \,   v_{n, k}^{-1}(\zeta)$, etc. We also denote $z_k\isdef \varphi_{n, k}(\zeta)$ and $z_{k+1}\isdef \varphi_{n, k+1}(\zeta)$. 

Substituting $z_k=\varphi_{n, k}(\zeta)$ in \eqref{eqCauchyTransfNew} and using that $ v_{n, k}'(z_k)=1/ \varphi_{n, k}'(\zeta)$, we obtain
$$
v_{n, k+1}\left( \varphi_{n, k}(\zeta) \right)
=\zeta +\frac{1}{n \zeta} \frac{1}{\varphi_{n, k}'(\zeta) }.
$$
Applying now $\varphi_{n, k+1}=v_{n, k+1}^{-1}$ to both sides yields
$$
\varphi_{n, k}(\zeta) =\varphi_{n, k+1}\left( \zeta +\frac{1}{n \zeta} \frac{1}{\varphi_{n, k}'(\zeta) }\right),
$$
and thus,
\begin{equation}
    \label{firstidentityPhi1}
\varphi_{n, k+1}(\zeta)- \varphi_{n, k}(\zeta) =  \varphi_{n, k+1}(\zeta) - \varphi_{n, k+1}\left( \zeta +\frac{1}{n \zeta} \frac{1}{\varphi_{n, k}'(\zeta) }\right).
\end{equation}
Using that for small $s\in \C$,
\begin{equation}
    \label{taylorEx}
 \varphi_{n, k+1}\left( \zeta +s\right)= \varphi_{n, k+1}\left( \zeta \right)+ \varphi_{n, k+1}'\left( \zeta \right) s + \frac{1}{2}\varphi_{n, k+1}''\left( \zeta +\theta \right) s^2,
\end{equation}
we get that 
\begin{equation}
    \label{identityPhi1}
    \begin{split}    
  z_{k+1}-z_k & =   \varphi_{n, k+1}(\zeta)- \varphi_{n, k}(\zeta)  =  -\frac{1}{n \zeta} \frac{\varphi_{n, k+1}'\left( \zeta \right)}{\varphi_{n, k}'(\zeta) }-\frac{\delta_{n, k}(\zeta)}{n^2 \zeta^2} \\
&  = -\frac{1}{n \zeta} -\frac{1}{n \zeta} \frac{\varphi_{n, k+1}'\left( \zeta \right)-\varphi_{n, k}'\left( \zeta \right)}{\varphi_{n, k}'(\zeta) }-\frac{\delta_{n, k}(\zeta)}{n^2 \zeta^2}. 
\end{split}
\end{equation}

With the notation introduced above, let us rewrite
\begin{equation} \label{identitiesderivativesPhi}
\frac{\varphi_{n, k+1}'\left( \zeta \right)-\varphi_{n, k}'(\zeta)}{\varphi_{n, k}'(\zeta)  }  = \frac{ v_{n, k}' \left( z_k \right) - v_{n, k+1}' \left( z_{k+1} \right) }{  v_{n,k+1}'\left( z_{k+1} \right)}  =  \frac{1}{n}\, \varepsilon_{n,k}^{(1)} (\zeta) - \varepsilon_{n,k}^{(2)} (\zeta)(u-z) ,
\end{equation}
where, again by \eqref{eqCauchyTransfNew}, 
\begin{equation}
\begin{split}
\varepsilon_{n,k}^{(1)} (\zeta) & = n \, \frac{  v_{n, k}'\left( z_{k+1}\right)-  v_{n, k+1}'(z_{k+1})}{ v'_{n, k+1}(z_{k+1}) } =- \frac{1}{   v'_{n, k+1}(z_{k+1}) } \,  \left(  \frac{  v'_{n,k}(u)}{  v_{n,k}(u)}\right)\bigg|_{u=z_{k+1}} \\
& =-  \frac{1}{z_{k+1}^2 \,  v'_{n, k+1}(z_{k+1}) } \, z_{k+1}^2 \left(  \frac{  v'_{n,k}(u)}{  v_{n,k}(u)}\right)'\bigg|_{u=z_{k+1}}
\label{expressionepsilonMod}
\end{split}
\end{equation}
and
$$
\varepsilon_{n,k}^{(2)} (\zeta) = \frac{ v_{n, k}' \left( z_k \right)- v_{n,k}'\left( z_{k+1} \right)}{    z_k-z_{k+1} } \, \frac{ 1}{  v_{n,k+1}'\left( z_{k+1} \right)  }.
$$
Combining \eqref{identityPhi1} and \eqref{identitiesderivativesPhi} we get that 
\begin{equation}
    \label{intermediate}
    \left(z_{k+1}-z_k\right) \left(1- \frac{\varepsilon_{n,k}^{(2)} (\zeta)}{n \zeta} \right) = -\frac{1}{n \zeta} -\frac{\varepsilon_{n,k}^{(1)} (\zeta)}{n^2 \zeta}  -\frac{\delta_{n, k}(\zeta)}{n^2 \zeta^2}.
\end{equation}
To prove \eqref{boundDelta}, it suffices to show that $\varepsilon_{n,k}^{(1)} (\zeta)$, $\varepsilon_{n,k}^{(2)} (\zeta)$, and $\delta_{n,k} (\zeta)$ are uniformly bounded for sufficiently small $|\zeta|$ (or equivalently sufficiently large $z_k$ and $z_{k+1}$). 

By \eqref{identitiesderivativesPhi} and Proposition~\ref{prob:boundsCauchy}, $\varepsilon_{n,k}^{(1)}(\zeta)$ is analytic for $|z_{k+1}|\ge R$, and functions
$$
u\frac{ v'_{n,k}(u)}{  v_{n,k}(u)} \quad \text{and} \quad \ u^2\,   v'_{n,k+1}(u) 
$$
are analytic and non-vanishing in $|u|\ge R$. By \eqref{boundsCauchy1}--\eqref{boundsCauchy2}, for $|u|\ge R$,
$$
 \left| u\, \frac{ v'_{n,k}(u)}{  v_{n,k}(u)} \right|= \left|  \frac{u^2    v'_{n,k}(u)}{ u\,  v_{n,k}(u)} \right|\le 18,
$$
and in consequence,
$$
\left| \frac{  v'_{n,k}(u)}{  v_{n,k}(u)} \right|\le \frac{18}{R},  \qquad |u|=R.
$$
By the Cauchy integral formula,
\begin{equation}
    \label{estimate1onPhi}
\left| u^2  \left( \frac{ v'_{n,k}(u)}{  v_{n,k}(u)} \right)'\right|\le 18 (R+1)^2,  \qquad |u|=R+1.
\end{equation}
On the other hand, by \eqref{boundsCauchy2},
\begin{equation}
    \label{boundPhi}
\left| \frac{1}{ u^2\,    v'_{n, k+1}(u) } \right| \le \frac{1}{1-(k+1)/n}\frac{(1+r/R)^4}{1-2r/R}\le \frac{4(3+2 \sqrt{2})}{1-(k+1)/n}, \quad |u|\ge R\ge (\sqrt{2}+1)r. 
\end{equation}
Using these estimates in \eqref{expressionepsilonMod}, we conclude that $\left| \varepsilon_{n,k}^{(1)}(\zeta)\right|$ is uniformly bounded for $|z_{k+1}|\ge R+1$ and $k/n\le T<1$. 

As for $\varepsilon_{n,k}^{(2)} (\zeta)$, we notice that for $z_k, z_{k+1}\in \Omega_R$,
$$
\left|  \frac{ v_{n, k}' \left( z_k \right)- v_{n,k}'\left( z_{k+1} \right)}{    z_k-z_{k+1} }\right|  \leq \max_{s\in \Omega_R} \left| v_{n,k}''\left( s \right)\right| .
$$
Appealing again to \eqref{boundsCauchy1} and Cauchy's formula, we get the uniform boundedness of $\varepsilon_{n,k}^{(2)} (\zeta)$ for  $z_k, z_{k+1}\in \Omega_{R+1}$.

Finally, by \eqref{taylorEx}, $\delta_{n, k}(\zeta)$ in \eqref{identityPhi1} satisfies   
\begin{equation}
    \label{eq:13}
\left|\delta_{n, k}(\zeta)\right| \leq \frac{1}{2 }   \max _{\zeta \in \D_\rho} \frac{ \varphi_{n, k+1}'' \left(\zeta\right)}{ \left( \varphi_{n, k}' \left(\zeta\right)\right)^2 },\quad \zeta\in \D_\rho.
\end{equation}
We have
\begin{align*}
\frac{ \varphi_{n, k+1}'' \left(\zeta\right)}{ \left( \varphi_{n, k}' \left(\zeta\right)\right)^2 } 
&= - \frac{   \,v_{n, k+1}''  \left(z_{k+1} \right)}{ \left(   v_{n, k+1}' \left(z_{k+1}\right)\right)^3 } \left(  v_{n, k}' \left(z_k\right)\right)^2 
= - \frac{   v''_{n, k+1} \left(z_{k+1} \right)}{    v'_{n, k+1} \left(z_{k+1}\right)  } \left( \frac{  v'_{n, k} \left(z_k\right)}{  v'_{n, k+1} \left(z_{k+1} \right)}\right)^2 
\\
& = - \frac{   v''_{n, k+1} \left(z_{k+1} \right)}{   v'_{n, k+1} \left(z_{k+1}\right)  } \left(  1+ \frac{\varepsilon_{n,k}(\zeta)}{n}\right)^2.
\end{align*}
Reasoning as above, by \eqref{boundsCauchy1},
$$
\left|   v''_{n, k+1} \left(u \right)\right| \le   \frac{4\left( 1-(k+1)/n\right)}{R}, \quad |u|=R+1,
$$
while, by \eqref{boundPhi},
$$
\left| \frac{1}{  v'_{n, k+1}(u) } \right| \le   \frac{4(3+2 \sqrt{2}) R^2}{1-(k+1)/n}, \quad |u|\ge R\ge (\sqrt{2}+1)r,
$$
showing that $\delta_{n,k}$ in \eqref{identityPhi1} is uniformly bounded also.
\end{proof}

\end{document}
